%======================================================================
%  Results -- Pretraining  (paper section 2; main figure = Figure 1)
%
%  Source artifact: GPN-Plasmids pretraining & model
%    https://claude.ai/code/artifact/07353969-1e26-4176-a170-db9d95886250
%    /home/canall/DNA_LM/plasmids_GPN/manuscript/artifact/page.html
%    snapshot 2026-08-28 (page.html mtime 2026-08-28 13:27)
%
%  This prose is the artifact's "Main text" band (seven paragraphs). NOTE:
%  manuscript/pretraining/01_results_sampling_design.md is a longer,
%  section-headed narrative of the same material and is NOT what was
%  transcribed here -- see CONVERSION_NOTES.md. The Methods files 02-05 and
%  the parameter table 06 ARE used, in supp_methods_02_pretraining.tex and
%  supp_tables.tex.
%
%  2026-09-01, author-directed: the seven transcribed paragraphs were
%  REGROUPED into three Results subsections (2.1 IMG/PR processing,
%  2.2 diversity-aware sampling, 2.3 pretraining and model architecture),
%  laid out 2 / 3 / 3. Two sentences changed paragraph -- old P1 s2 opens
%  2.3, and old P2 s1+s3 open 2.2 -- and three sentences were added: the
%  Fig. 1c callout in 2.1, the closer to 2.1, and the second half of the
%  2.2 opener. No transcribed clause was deleted. See CONVERSION_NOTES.md.
%
%  Superscript citations 1-9 in the artifact map to the merged bibliography:
%    1 camargo2024img  2 devlin2018bert  3 benegas2023dna
%    4 redondosalvo2020pathways  5 jain2018high  6 nguyen2024sequence
%    7 kalchbrenner2016neural  8 mukherjee2023twenty  9 shaw2023fast
%======================================================================

% ---- 2.1  IMG/PR processing ------------------------------------------
% \subsection*{A quality-filtered plasmid corpus resolved into a cluster hierarchy}
\subsection*{IMG/PR processing}
\label{sec:imgpr}

\model is a genomic language model trained on 693,638 plasmid sequences drawn
from the Integrated Microbial Genomes / Plasmid Repository (IMG/PR) collection
\cite{camargo2024img}, after quality filtering and a PTU-level species holdout
(\nameref{sec:methods}). Host assignments in the corpus span 48 phyla, and its
ecosystem composition, annotated through GOLD, \cite{mukherjee2023twenty} is
shown alongside that host phylogeny in \fref{fig:pretraining}c. In contrast to
genomic language models trained across broad evolutionary timescales, \model is
trained exclusively on mobile genetic elements, without an intentionally
included chromosomal reference corpus.

Plasmids have no universal hierarchical classification. The available grouping,
the Plasmid Taxonomic Unit (PTU), \cite{redondosalvo2020pathways} is a broad
partition of undefined phylogenetic rank that says nothing about structure
within the group it names. To expose the missing structure we clustered every
sequence at two levels --- dereplication clusters at $\geq$99\% average
nucleotide identity (ANI) with $\geq$95\% bidirectional alignment fraction, and
sub-PTUs at $\geq$95\% ANI \cite{jain2018high} with $\geq$95\% unidirectional
alignment fraction --- which nest exactly inside PTUs
(\edfref{fig:ed-corpus}f,g). Classifying each PTU by which levels it populates
reveals a sharp asymmetry: singleton (Unique) PTUs, holding a single
dereplication cluster, are 69\% of all PTUs but only 30\% of dereplicated
sequences, whereas the 19\% of PTUs that resolve into multiple sub-PTUs (Major)
hold 59\% (\fref{fig:pretraining}d). That asymmetry is what any sampling policy
over this corpus must contend with.


% ---- 2.2  Diversity-aware sampling for pretraining --------------------
% \subsection*{A diversity-aware epoch policy between two unsatisfactory extremes}
\subsection*{Diversity-aware sampling for pretraining}
\label{sec:sampling}

Training a language model on this corpus raises a sampling problem that does
not arise for cellular genomes. With no taxonomy to sample across, abundance in
a plasmid database reflects sequencing effort rather than diversity, so a
training budget spent in proportion to abundance is spent in proportion to how
often a clade has been sequenced.

Near-identical sequences are first collapsed by dereplication, so no plasmid
sequence is over-weighted merely for having been resequenced; even so, the
two obvious sampling policies are unsatisfactory extremes
(\fref{fig:pretraining}e). After dereplication, sampling every remaining
sequence with equal probability concentrates the budget on the largest clades
--- the most diverse PTU is drawn 715$\times$ more often than under an even
one-sequence-per-PTU budget, and the top 1\% of PTUs absorb 24\% of all
sampled sequences --- so the model is optimized largely on a handful of
well-sampled clades. Sampling exactly one sequence per PTU, the strategy Evo
applies to IMG/PR, \cite{nguyen2024sequence} removes that imbalance but
discards all within-clade variation and carries a robustness cost specific to
a metagenomic corpus. Plasmid identity in IMG/PR is assigned computationally,
and a misclassified chromosomal fragment has no genuine plasmid sequence to
group with, so it tends to survive clustering as its own singleton. Because
Unique PTUs are 69\% of all PTUs, a one-sequence-per-PTU budget would spend
the majority of training on single, uncorroborated sequences.

We therefore constructed each training epoch by a diversity-aware policy that
sits deliberately between these extremes (\fref{fig:pretraining}f,
\nameref{sec:methods}). For Major PTUs we implemented a sampling scheme that
scales sub-linearly with the size of the individual Major PTU, measured by
its number of sub-PTU clusters $n_i$: the per-PTU allocation grows as
$\min(n_i, Q_{0.95})^{1/T}$, where the cap $Q_{0.95}$ at the 95th percentile
of sub-PTU count controls for the sequencing bias inherent in metagenomic
databases such as IMG/PR, and the temperature $T = 1.5$ sets the sub-linear
scaling, so that additional sequences from an already well-sampled clade
yield diminishing returns (\edfref{fig:ed-corpus}i,
\nameref{sec:methods}). A floor of one guarantees every Major PTU a
representative, so within-clade variation is retained without any clade
dominating; Moderate PTUs contribute exactly one sequence, and Unique PTUs
are uniformly subsampled each epoch, with 56.4\% drawn. The cap fixes an
expected-allocation ceiling of 3.15 representatives per epoch for any PTU
regardless of size --- including the PTU containing 735 sub-PTUs.
% Next sentence, 2026-09-01, author-directed: the epoch size read
% "approximately 170,000" in the transcribed artifact. Reconciled to the
% measured value, which matches supp_methods_02_pretraining.tex (mean total
% 166,192, s.d. 2.9) and the Fig. 1d caption. See CONVERSION_NOTES.md.
The resulting epoch, a mean of 166,192 sequence representatives, is
partitioned 35\% / 15\% / 50\% across the three categories: this pulls the
most over-represented PTU from 715$\times$ its even per-PTU share to
4.0$\times$ and the top 1\% of PTUs from 24\% of sampled sequences to 4\%
(\fref{fig:pretraining}e), redistributing budget toward internally diverse
Major clades while reserving half for the Unique majority. The whole
procedure is re-executed at every epoch boundary, so representatives,
allocations, window phases and masks are redrawn across approximately 102
epochs of training.


% ---- 2.3  Pretraining and model architecture --------------------------
% \subsection*{A 468M-parameter dilated convolutional encoder trained by masked prediction}
\subsection*{Pretraining and model architecture}
\label{sec:pretraining}

\model is trained by masked language modelling \cite{devlin2018bert} in a
convolutional framework adapted from GPN \cite{benegas2023dna} to learn a
nucleotide probability distribution conditioned on surrounding sequence
(\fref{fig:pretraining}g). The model is a 467.9M-parameter encoder of 72
residual blocks of dilated, non-causal convolutions adapted from ByteNet,
\cite{kalchbrenner2016neural} with a hidden dimension of 768. Dilation cycles
through 1 to 32 every six layers, giving a receptive field of 6,049 bp ---
roughly three times the 2,048 bp input window, so every output position
conditions on the whole window without attention (\fref{fig:pretraining}g).
Nucleotides enter as one-hot vectors with no learned embedding and no
positional encoding, since the convolutional stack is translation-equivariant
by construction. Complete plasmids are circularized before windowing so that
predictions near a window edge are informed by genuine flanking sequence
rather than an artefact of linearization (\suppfref{fig:supp-windowing}).
Full architectural and training details are provided in
\nameref{sec:methods}.

\model was trained from random initialization for 375,000 optimizer steps
under a constant learning rate with warmup, consuming approximately 217
million sequence windows or $4.5\times10^{11}$ nucleotide tokens in BF16
across up to eight NVIDIA H200 GPUs.

The purpose of modelling plasmid sequence space in this way is discovery in
the parts of it that are least well characterized. Origins of transfer and
replication, and the accessory modules that accompany them, are best
described in a handful of clinically important families and remain largely
uncharacterized in the metagenomic majority. We next evaluate whether these
representations identify origins of transfer and replication across diverse
plasmid sequences.
